\documentclass[letterpaper,11pt,preprint]{aastex}
\usepackage{graphics,graphicx}
\usepackage{amsmath}
\usepackage{natbib}
\usepackage{color}
\citestyle{aas}

\setlength{\textwidth}{6.5in} \setlength{\textheight}{9in}
\setlength{\topmargin}{-0.0625in} \setlength{\oddsidemargin}{0in}
\setlength{\evensidemargin}{0in} \setlength{\headheight}{0in}
\setlength{\headsep}{0in} \setlength{\hoffset}{0in}
\setlength{\voffset}{0in}

\makeatletter
\renewcommand{\section}{\@startsection%                                                                                                                                                                    
{section}{1}{0mm}{-\baselineskip}%                                                                                                                                                                         
{0.5\baselineskip}{\normalfont\Large\bfseries}}%                                                                                                                                                           
\makeatother

%%%%%%%%%%%%%%%%%%%%%%%%%%%%%                                                                                                                                                                              
%%%%% Start of document %%%%%                                                                                                                                                                              
%%%%%%%%%%%%%%%%%%%%%%%%%%%%%                                                                                                                                                                              

\begin{document}
\pagestyle{plain}

\section{How Loud are Sounds?}

In the previous note, we saw that atmsopheric pressure is close to
100,000 Pascals, equivalent to about 1 kg per square centimetre.  It
may not surprise you to hear that what we hear as sound comes from
changing pressure, so when we talk about how ``loud'' a sound is, what
we really mean from a physics perspective is ``how much is the
pressure changing?''  It's probably not a lot, at least relative to
atmospheric pressure, because sounds don't knock us over while a good
gust of wind would.  How much lower, though?

We'll try to make a reasonable estimate based on some sensible
guesses.  The answer won't be perfect, but we'll at least be in a
position to say ``is it a thousandth of atmospheric pressure?  A
millionth?  A billionth?''  This sort of answer is called an
``order-of-magnitude estimate'' -- getting a perfect answer would be
much, much harder, but at least we can get started with our
order-of-magnitude guess.  Instruments like flutes or violins or
trombones would be hard to estimate, because it's really difficult to
know how much of the work a musician is doing ends up as sound.  In
fact, the answer is usually ``not very much'' because most of the
energy ends up getting dissipated as heat in the instrument rather
than coming out as sound.  One case we can estimate is a piano.
Because of the way a piano is played, we can put an upper limit on how
much energy goes into sound.  Presumably nearly all of you have heard
a piano in a concert at some point as well.  

\subsection{How Much Work in a Piano Note?}
The most basic definition of work is ``force times distance.''  If I
know how much force I exerted on something, and how far I pushed it, I
know how much work I did.  The metric unit of work or energy is a
Joule, the work required to apply a force of one Newton over one
metre.  As an example, how much work does it take to climb a flight of
stairs?  Let's say you're 100 kg, and a flight of stairs goes up 3m.
Newton's second law is $F=ma$, force equals mass times acceleration,
and at the Earth's surface, the acceleration from gravity is about 10
$m/s^2$.  The force from gravity on us is then $100 kg \times 10 m/s^2
= 1000 kg m/s^2=1000 N$.  To climb the stairs, we have to exert a
force of at least 1000 N upwards to counteract gravity, and that force
is doing the work.  If we climb 3$m$, then force times distance is
$1000N \times 3m = 3000 N-m = 3000J = 3kJ$.  In most of the world, the
energy in food is listed directly in $kJ$, but here in Canada, it's
usually in kilo-calories, which we usually just call calories.  One
(kilo-)calorie is 4.19 $kJ$, so it takes a bit less than a calorie to
climb a flight of stairs.

Let's redo the same calculation, except now for a piano.  We don't
know how much of the work we do will end up as sound, but we know it
can't be more than the work we did pressing a key.  How much work do
we do pressing a key?  Work is force times distance, so let's start
with a distance.  Picture a piano in your head - about how far does a
key go down?  It's a lot more than a millimetre, but a lot less than a
metre, so 1 cm seems reasonable.  Your intrepid instructor measured
the key travel on the piano at his house, and indeed the travel is
almost exactly 1 cm.  How about force?  One way to estimate is to
imagine different weights on a piano key.  What sort of weight would
it take resting on a key to push the key down and make a sound?  Well,
your instructor set a bottle of whiskey on the keyboard and let it
push some keys down.  The bottle weight 1.3 kg, and it pushed down 4
keys at what sounded like a \textit{mezzo-piano} to my ear.  The total
force is $1.3*10=13 N$ split amongst 4 keys, so the force on each key
was $13/4 \sim 3N$.  A force of $3N$ over a distance of $1cm=0.01m$
then is a total work of $0.03J$.  The amount of energy might be (much)
less than this, but it cannot be more.

Since we can hear a pianist playing at a \textit{mezzo-piano} volume
inside a concert hall, we know that $0.03J$ will be easily heard.  How
do we convert that work to a change in pressure, though?  For
simplicity, we can assume we have a partition in our concert hall with
half the air on one side and half on the other.  We push on that
partition, compressing the air on one side and expanding it on the
other, until we've done our $0.03J$ of work.  You're hopefully now
getting an intuitive sense for just how small the pressure changes
from typical sounds are.  A thirtieth of a Joule, spread across a
concert hall, is not going to get you a very large pressure change.

Look at a picture of Montreal's Maison Symphonique.  The stage is
about 20m across, the hall is about as high as it is wide, and it's a
few times longer than it is wide.  For ease, let's call the dimensions
20x20x50m.  I set up my partition that splits the hall in half, so now
I have two volumes of 10x20x50, and I'm going to move along the 10m
direction.  What is the net force on the partition?  Well, if I
haven't moved the partition yet, the net force is zero because the
pressures on both sides match.  What happens after I move it a
distance $dx$ though?  Well, now the density on one side has gone up
by a factor of $dx/10$, and the density on the other side has gone
down by the same amount.  For simplicity, we'll assume the temperature
stays the same\footnote{In fact, this is not a good assumption as the
air on the compressed side heats up as well, but we can ignore that
for our order-of-magnitude estimate.}.  From the ideal gas law with
temperature fixed, we know pressure times volume is a constant, so
whatever factor the volume went down is, the pressure must have gone
up by the same factor.  Let's say we  move our partition by 1 cm, then
what happens?  The new volume is (10m-0.01m)/(10m) = (1-0.001)m, or
one part in a thousand smaller.  That means the new pressure is a part
in a thousand larger, so the change in pressure is 100 kPa/1000=100
Pascals.  Force is pressure times area, so the force on our partition
100 times the area of the partition.  We're moving in the 10m
direction, so the area is 20m by 50m, or $1000m^2$.  The total force
on the partition is 100*1000 or 100,000N.  The work we did was that
force times the 1cm distance we moved, so $100,000N \times 0.01m =
1000J$.  That's way, way more than the 0.03J we have, so the distance
moved must be much less than 1 cm.  At this point, the easiest way to
get our actual answer is to take what happens at 1 cm, and adjust it
to get the answer we want.  If we move by a distance $dx$, then the
pressure changes by a factor of $dx/x$, so the force must be
proportial to $dx/x$ as well.  The work is force times distance, and
the distance is $dx$, so the work must scale like $dx^2/x$.  Since $x$
doesn't change (after all, the room size is fixed), so the work we do
goes like $dx^2$.  The actual work we did was 0.03J, so we must have
\begin{eqnarray}
  (dx/0.01)^2 = (0.03/1000) \\
  dx/0.01 = \sqrt{0.03/1000}\\
  dx=0.01\times \sqrt{0.03/1000}=5.47\times 10^{-5}m \sim 50 \mu m
\end{eqnarray}


Now we have at least a rough estimate for sound pressure.  Our piano
note in a concert hall moves the partition by about 50 microns out of
10 metres, a fractional change of $5\times 10^{-6}$.  A fractional
volume change of five parts per million will lead to a similar
fractional pressure change, and since the pressure is 100,000 Pascals,
the \textit{change} in pressure is on the order of half a Pascal.  We
know this is an overestimate, because we assumed all of the energy
from pressing the piano key ended up as sound, and that none of that
sound leaked out of the concert hall, but at least we're in the
ballpark.  If you go to \textit{e.g.} Wikipedia, you'll see that sound
pressures of half a Pascal is actually really loud - about 88 dB, or
just above the threshold where hearing protection is required to
prevent hearing loss.  The true answer is about a factor of 10 or so
lower, or an order-of-magnitude.  We call these order-of-magnitude
estimates for a reason!

For fun, let's see if we can do even better.  There are two main
assumptions we made that probably weren't very good.  The first is
that all the energy in the key press ended up as sound.  It's tricky
to estimate this well - to get an idea, there's a paper by Seok Jae
Bang (2009) that estimates the acoustic efficiency of a Yamaha
upright.  They find that the efficiency is only about 1\%, so we
should decrease our sound energy by a factor of 100.  Second, we
assumed that all the sound energy ended up in the hall \textit{at the
  same time}.  That's also obviously wrong.  Why?  Well, if I play a
piano note and hold the key down, it takes quite a while (many
seconds) for the sound volume to drop, so the energy takes a while to
go from vibrations of the piano strings to vibrations of the air, say
10 seconds.  However, in a concert hall, if I make a noise, the time
it takes the sound volume to drop (the reverb time) is a lot less,
more like a second.  That means that at any given time, only about 1s
(reverb)/10s (piano sustain) or 10\% of the total sound energy we'll
put out is contained in the hall.  With these correction factors, our
sound energy is down by a factor of 1000 from our initial
overestimate, so the pressure is down by a factor of the square root
of that, so a better estimate is 0.01 Pascal for our sound pressure.
That's actually pretty good - Wikipedia tells us that's in the range
of normal conversation (0.002-0.02 Pascals, or 40-60 dB).

One final calculation, also for fun, I'd like to do is to compare the
sound pressure we just calculated to change in pressure from changing
altitude.  The height scale of the atmosphere is close to 10 km, which
is why when you climb a tall mountain you notice there's a lot less
air, but not when you climb stairs.  If we assume the pressure goes
from 100 kPa to 0 when we go up 10 km (again, we're after a rough
estimate), then the change per metre is 100,000/(10km=10,000m), or 10
Pascals/metre.  Just standing up will change the pressure by 10
Pascals, which is much, much bigger than the sounds!  Our 0.01 Pascal
improved estimate for a piano volume corresponds to 0.01/10=0.001m or 1
mm.  Because we (usually) change altitude slowly and smoothly, our ears
do not register those pressure changes as sounds.  If they did, we'd
all go deaf just from getting out of bed in the morning.  




\end{document}
